% ═════════════════════════════════════════════════════════════════════════════
\section{Parameterized Systems}
\label{sec:prelim}
% ═════════════════════════════════════════════════════════════════════════════

This section fixes the setting for the rest of the chapter. We describe
parameterized systems and the verification problem informally, introduce a
running example, and recall the basic notions of well-quasi-orderings and
monotonic transition systems, which underlie several of the techniques
presented later. We keep the presentation informal; precise definitions for
the different classes of systems can be found in the cited
works~\cite{AST18,BJKKRVW15,Abd10}.

\subsection{Process Templates and Configurations}

A parameterized system is defined by a \emph{process template}, a finite-state
machine that describes the behavior of a single process. The template has a
finite set~$Q$ of \emph{local states}, a set of initial local states, and a set
of transitions. For each natural number~$n$, the \emph{instance} of size~$n$
consists of~$n$ copies of the template running in parallel, where the copies
are identified by the indices $1, \ldots, n$. Some systems also contain a
distinguished \emph{control process}, defined by a separate template, that runs
in parallel with the copies.

A \emph{configuration} of the instance of size~$n$ records the local state of
each process. We represent it as a word $q_1 q_2 \cdots q_n$ of length~$n$
over the alphabet~$Q$, where $q_i$ is the local state of process~$i$. When the
system also has shared variables, the configuration includes their values. For
each fixed~$n$, the instance is a finite transition system. The parameterized
system is the infinite family of all instances, and its set of configurations,
taken over all~$n$, is infinite.

The \emph{topology} of a system determines how processes are arranged and
which processes can observe each other. In a \emph{linear} topology, the
processes form an array, and a process can distinguish the processes to its
left from those to its right. In a \emph{ring}, the last process is connected
back to the first. In a \emph{clique}, all processes are indistinguishable; the
order of the letters in a configuration is then irrelevant, and a configuration
can be represented by a multiset of local states, or equivalently by a vector
that counts, for each local state, how many processes occupy it. Trees and
more general graphs also occur, for instance in hierarchical protocols and in
ad hoc networks.

Processes interact through a small number of communication primitives. A
\emph{local} transition changes the state of a single process without
involving the others. In a \emph{rendez-vous}, two processes synchronize and
change their states simultaneously. In a \emph{broadcast}, one process sends a
message and all other processes react to it in the same step. Processes may
also read and write \emph{shared variables} with finite domains. Finally, a
transition may be guarded by a \emph{global condition} that refers to the
states of the other processes. The condition is \emph{universal} if it requires
all other processes (or all processes to the left, or to the right) to be in
given local states, and \emph{existential} if it requires at least one of them
to be. Global conditions are a convenient way of modeling protocols in which a
process inspects the states of its peers before taking a step.

\subsection{The Verification Problem}

Most of this chapter is concerned with safety properties. A safety property is
given by a set~$\Bad$ of \emph{bad} configurations, which is defined
independently of the size of the instance. For instance, for mutual exclusion,
$\Bad$ is the set of configurations, of any length, in which at least two
processes are in their critical sections. We write $\Init$ for the set of
initial configurations of all instances. The \emph{parameterized reachability
problem} asks whether there is a size~$n$ and a run of the instance of
size~$n$ that starts from a configuration in~$\Init$ and reaches a
configuration in~$\Bad$. The system is \emph{safe} if no such~$n$ and run
exist, that is, if the union of the reachable sets of all instances is
disjoint from~$\Bad$.

More general properties are often expressed in \emph{indexed} temporal
logics~\cite{BCG89,EN95,GerSis92}, whose formulas quantify over process
indices. A formula of the form ``for all distinct indices~$i$ and~$j$,
$\varphi(i,j)$'' states that the temporal property~$\varphi$ holds for every
pair of processes. Liveness properties, such as the requirement that every
process that requests access to its critical section eventually obtains it,
are typically checked under fairness assumptions on the scheduler.

As mentioned in Section~\ref{sec:intro}, parameterized verification is
undecidable in general~\cite{AK86}. The proof encodes the tape of a Turing machine in the configuration of a linear array of
processes, one process per tape cell, with the head stored in the local state
of one of the processes. Since the number of processes is unbounded, every
finite portion of the tape can be represented, and the halting problem reduces
to the parameterized reachability problem. The same idea works for rings of
identical processes~\cite{Suz88}. Each technique in this chapter can be
seen as a response to this result: it either identifies a class of systems for
which the problem is decidable, or accepts incompleteness in exchange for
generality.

\subsection{A Running Example}
\label{sec:running}

We use the following simple mutual exclusion protocol as a running example.
Each process has three local states: $\Idle$ (idle), $\Wait$ (waiting), and
$\Crit$ (in the critical section). Initially, all processes are in~$\Idle$. A
process may move from~$\Idle$ to~$\Wait$ at any time. A process in~$\Wait$ may
enter the critical section, moving to~$\Crit$, if no other process is
in~$\Crit$; this is a universal global condition. A process in~$\Crit$ may
return to~$\Idle$ at any time. The template is shown in
Figure~\ref{fig:running}.

\begin{figure}[t]
\centering
\begin{tikzpicture}[->,>={Stealth[round]},shorten >=1pt,auto,
    node distance=2.2cm,semithick,
    every state/.style={minimum size=8mm,inner sep=1pt}]
  \node[state,initial,initial text={}] (i) {$\Idle$};
  \node[state] (w) [right=of i] {$\Wait$};
  \node[state] (c) [right=3.6cm of w] {$\Crit$};
  \path (i) edge node {} (w)
        (w) edge node[above=1pt] {\small $\forall j \neq \mathit{self}\colon q_j \neq \Crit$} (c)
        (c) edge[bend left=35] node {} (i);
\end{tikzpicture}
\caption{The process template of the running example. The transition from
$\Wait$ to~$\Crit$ is guarded by a universal condition: it is enabled only if
no other process is in~$\Crit$.}
\label{fig:running}
\end{figure}

Configurations are words over the alphabet $\{\Idle, \Wait, \Crit\}$, and the
initial configurations are the words $\Idle^n$, for $n \geq 1$. The bad
configurations are the words that contain at least two occurrences
of~$\Crit$. For example, the instance of size three can reach
$\Wait\Idle\Wait$ from $\Idle\Idle\Idle$ in two steps, and then move to
$\Crit\Idle\Wait$ and to $\Idle\Idle\Wait$, but it can never reach a
configuration such as $\Crit\Idle\Crit$, since the second process in~$\Wait$
is blocked as long as the first one is in~$\Crit$.

The protocol is safe for every~$n$, and for a human reader this is obvious:
the guard is checked and the state is changed in a single atomic step, so no
process can enter the critical section while another one is there. The
example is therefore not interesting as a verification challenge, but it is
small enough to show how each technique handles the unbounded number of
processes. Realistic protocols, such as Burns' or Szymanski's algorithms,
replace the atomic universal condition by a loop in which a process inspects
the other processes one at a time. Their correctness then depends on subtle
interactions between interleaved inspections, and verifying them is
considerably harder~\cite{ABDR08,AHH16}.

\subsection{Well-Quasi-Orderings and Monotonic Systems}
\label{sec:wqo}

Several techniques in this chapter rely on orderings on configurations that
compare instances of different sizes. A \emph{quasi-ordering} on a set is a
reflexive and transitive relation~$\sqsubseteq$. We say that $\sqsubseteq$ is a
\emph{well-quasi-ordering} (WQO) if every infinite sequence $a_0, a_1, a_2,
\ldots$ contains two elements $a_i$ and $a_j$ with $i < j$ and $a_i \sqsubseteq
a_j$. For words, we use the \emph{subword ordering}. We write $u \preceq v$ if
$u$ can be obtained from~$v$ by deleting letters, that is, if $u$ occurs in~$v$
as a not-necessarily-contiguous subsequence. For example, $\Wait\Crit \preceq
\Wait\Idle\Idle\Crit$. Two classical results provide the WQOs used most often
in parameterized verification. By Dickson's lemma, vectors of natural numbers
of a fixed dimension, ordered componentwise, form a WQO. By Higman's
lemma~\cite{Hig52}, the words over a finite alphabet form a WQO under the
subword ordering. Higman's lemma extends to words over any
well-quasi-ordered alphabet, which makes it possible to combine orderings in a
modular way.

A set~$U$ of configurations is \emph{upward-closed} if it contains, together
with every element~$a$, every element~$b$ with $a \sqsubseteq b$. Symmetrically,
$U$ is \emph{downward-closed} if it contains, together with every element~$b$,
every element~$a$ with $a \sqsubseteq b$. We write $\upc{a}$ for the set of elements larger than or equal to~$a$. In the running
example, the set of bad configurations is upward-closed under the subword
ordering: it is the set $\upc{(\Crit\Crit)}$ of words that contain
$\Crit\Crit$ as a subword. Upward-closed sets play a central role for two
reasons. First, under a WQO, every upward-closed set is determined by its
finitely many minimal elements, which gives a finite representation of an
infinite set. Second, every increasing sequence $U_0 \subseteq U_1 \subseteq
\cdots$ of upward-closed sets eventually stabilizes.

We say that a transition system is \emph{monotonic} with respect to an
ordering~$\sqsubseteq$ on its configurations if larger configurations can
simulate smaller ones: whenever $a$ can move to~$a'$ in one step and $a
\sqsubseteq b$, the configuration~$b$ can move in one step to a
configuration~$b'$ with $a' \sqsubseteq b'$. A \emph{well-structured transition system} is a
transition system that is monotonic with respect to a
WQO~\cite{ACJT96,ACJT00,FS01}. In a well-structured system, the set of
configurations from which an upward-closed set~$U$ can be reached is itself
upward-closed. This set can be computed by \emph{backward reachability}: one
starts from $U_0 = U$ and computes $U_{i+1} = U_i \cup \Pre(U_i)$, where
$\Pre(U_i)$ is the set of immediate predecessors of~$U_i$. Each $U_i$ is
upward-closed and is represented by its minimal elements, and the sequence
stabilizes after finitely many steps. Safety then amounts to checking that
no initial configuration belongs to the limit. The algorithm requires that the
minimal elements of $\Pre(U)$ can be computed from those of~$U$, which is the
case for most models considered in this chapter.

Many parameterized systems are monotonic under the subword ordering. Local
transitions, rendez-vous, broadcast, and existential global conditions do not
depend on the number of other processes in a way that could disable them, so
adding processes preserves the ability to perform them. Universal global
conditions are different. In the running example, the configuration~$\Wait$
can move to~$\Crit$, but the larger configuration $\Wait\Crit$ cannot move its
waiting process, because the added process is in the critical section. The
running example is therefore not monotonic, and neither are most mutual
exclusion protocols. Sections~\ref{sec:cutoff} and~\ref{sec:rmc} describe two
ways of dealing with this: view abstraction, which does not require
monotonicity, and monotonic abstraction, which restores it by
over-approximation.

For systems with a clique topology, a configuration can be represented by a
vector of counters, one for each local state. This \emph{counting abstraction}
is exact when processes are anonymous, and it turns the system into a Petri
net, or equivalently a vector addition system, in which each transition adds a
fixed vector to the current counter values~\cite{GerSis92,Esp98}. Petri nets
are monotonic with respect to the componentwise ordering, which is a WQO by
Dickson's lemma, so their coverability problem, and hence the parameterized
reachability problem for such systems, is decidable. A universal global
condition corresponds to a test that a counter is zero, which Petri nets
cannot express. This is another view of why universal conditions are the main
source of difficulty in the analysis of parameterized systems.
